\section{Audio Transformers: Raw Waveform y representaciones multiescala}

El tercer trabajo ya no predice el siguiente sample, sino que reconoce sound events a partir de un segundo de raw waveform. Retoma la idea central de ViT: dividir la entrada en patches, mapearlos a embeddings mediante un learned front end y después agregarlos con un Transformer; además, usa pooling o wavelet para que las capas posteriores vean un rango temporal más amplio. Aquí, «raw» todavía incluye la división en bloques y la proyección lineal del front end; no significa aplicar global attention a los 16,000 samples sin modificación alguna.

\lecturefigure{slide-38-audio-transformers-title.jpg}{Tercer trabajo: large-scale raw-audio understanding sin convoluciones.}{Video oficial de Stanford Online 00:37:50--00:38:21; Audio Transformers v1, arXiv:2105.00335v1.}

\subsection{La tarea FSD50K y los límites de una comparación justa}

En clase se usa FSD50K: unos cincuenta mil sound events human-labeled y unas 200 clases. La configuración de clase que muestra la slide es de unos 40,000 training clips y 10,000 validation/test clips, alrededor de cien horas en total; la entrada se submuestrea a 16 kHz y la muestra básica es una waveform de un segundo. El objetivo es multi-label audio content classification, no waveform reconstruction.

\lecturefigure{slide-39-fsd50k-dataset-setup.jpg}{FSD50K y configuración experimental: raw waveform, duración fija, multi-label content prediction.}{Video oficial de Stanford Online 00:37:54--00:38:42.}

Para una muestra $x$ y $C$ etiquetas, el modelo produce logits $s_c$; lo habitual es usar binary cross-entropy:
\[
\mathcal{L}_{\mathrm{BCE}}
=-\sum_{c=1}^{C}
\left[y_c\log\sigma(s_c)+(1-y_c)\log(1-\sigma(s_c))\right].
\]
donde $y_c\in\{0,1\}$ indica si la clase $c$ está presente. La evaluación usa mean average precision (mAP), que integra el precision--recall ranking de cada clase y después promedia sobre las clases; por eso no puede compararse directamente con la top-5 next-sample accuracy de las secciones anteriores.

\begin{warningbox}{Ambas se llaman Accuracy, pero las tareas son completamente distintas}
La raw synthesis anterior hace single-label next-step prediction entre 256 sample states; aquí, FSD50K es multi-label event detection y se reporta mAP. Las magnitudes numéricas de ambas no son comparables: no puede decirse que 0.537 sea «menor» que 85\%.
\end{warningbox}

\subsection{Por qué hacer Multi-Scale sobre el Intermediate Embedding}

La raw waveform contiene estructura a múltiples escalas, desde los transitorios hasta los eventos. Si solo se apilan Transformer layers de la misma longitud, aunque en teoría cada capa puede direccionar globalmente, la dimensión de la representación y la token rate no forman una jerarquía por sí solas. El artículo introduce pooling o wavelet para reducir la resolución temporal entre algunas capas, de modo que las capas superiores agreguen un rango temporal más largo con menos positions.

\lecturefigure{slide-40-wavelet-motivation.jpg}{Wavelet motivation: construir explícitamente una jerarquía multiescala sobre los intermediate embeddings.}{Video oficial de Stanford Online 00:38:43--00:39:28.}

\begin{knowledgebox}{Concepto previo: Receptive field y Resolution}
Receptive field es el rango temporal de la entrada del que puede depender una representación de alto nivel; resolution es cuántas posiciones conserva todavía el eje temporal. El pooling amplía lo que cubre cada posición, pero pierde el detalle fino; la wavelet produce a la vez una approximation de baja frecuencia y un detail de alta frecuencia, por lo que en teoría conserva más variación que un promedio simple, aunque la forma de combinarlos sigue determinando si la información realmente se aprovecha.
\end{knowledgebox}

\teachervoice{El ponente describe la wavelet como «aprender una hierarchy en forma de árbol dentro de una capa». Esto no significa que el modelo aprenda los parámetros de la wavelet; en clase se usa una Haar-like operation fija, y el aprendizaje ocurre en los Transformer embeddings y en las capas posteriores.}

\subsection{Haar Wavelet: promedio y diferencia}

Para embeddings adyacentes $h_{2i},h_{2i+1}\in\mathbb{R}^d$, la Haar transform más simple puede escribirse como
\[
a_i=\frac{h_{2i}+h_{2i+1}}{\sqrt{2}},
\qquad
d_i=\frac{h_{2i}-h_{2i+1}}{\sqrt{2}}.
\]
$a_i$ es la approximation de baja frecuencia, que resume el promedio local; $d_i$ es el detail, que conserva la variación entre vecinos. Si el sistema solo conserva el promedio, equivale a un pooling de submuestreo; si conserva o combina el detail, puede transportar las diferencias transitorias. Cuando la longitud pasa de $L$ a aproximadamente $L/2$, el número de attention pairs posteriores baja de $L^2$ a aproximadamente $L^2/4$.

\lecturefigure{slide-41-wavelet-in-transformer.jpg}{Inserción de una Haar wavelet/average pathway entre Transformer layers.}{Video oficial de Stanford Online 00:39:29--00:41:18.}

\begin{importantbox}{Por qué «sin parámetros nuevos» aún cambia el modelo}
La Haar transform no tiene learned weights, pero cambia la sequence length, la respuesta en frecuencia, el attention graph posterior y el gradient path. Que el número de parámetros no cambie no implica que no cambien el cómputo, la información ni el inductive bias.
\end{importantbox}

\subsection{Architecture: Front End, Transformer y Temporal Reduction}

El modelo primero introduce los waveform patches en un learned feed-forward front end, que los comprime en embeddings de menor dimensión; después pasan por varias Transformer layers, con average pooling o wavelet pathway entre ellas; al final se agregan y se producen los scores multietiqueta. El large model del artículo usa 6 layers, 64-dimensional embeddings, 8 heads y feed-forward blocks, con unos 2.3M parámetros; lo importante es comparar la eficiencia de representación con los baselines convolucionales de la misma tabla.

\lecturefigure{slide-42-audio-transformer-architecture.jpg}{Audio Transformer sobre raw waveform: front-end patches, attention layers, temporal reduction y classifier.}{Video oficial de Stanford Online 00:41:19--00:42:21.}

\lecturefigure{slide-43-learned-front-end.jpg}{Methodology: el front end aprende la representación tiempo-frecuencia, el Transformer aprende el embedding y pooling/wavelet amplían la escala.}{Video oficial de Stanford Online 00:42:22--00:44:01.}

\begin{knowledgebox}{Lectura de la figura: no confundir el Front End con «sin convoluciones no hay preprocesamiento»}
La parte izquierda todavía organiza la waveform en patches y usa un 2048-unit front end y una 64-dimensional projection para generar tokens; lo que recibe el Transformer son learned embeddings. Cuando el artículo dice que no hay convolution, se refiere a que no hay un convolutional backbone tradicional, no a que los samples de entrada no pasen por ningún mapeo local.
\end{knowledgebox}

\subsection{Tabla de resultados: eficiencia de parámetros y alcance de las métricas}

La tabla de la slide y del artículo muestra: VGG-like obtiene 0.434 mAP; el small Transformer, 0.469; el large 6-layer Transformer, 0.525, y el modelo del mismo tamaño con pooling, 0.537. CRNN, ResNet-18 y DenseNet-121 obtienen en esa tabla 0.417, 0.373 y 0.425, respectivamente. Los resultados respaldan que el raw-waveform Transformer es competitivo en esta configuración de FSD50K y también sugieren que la temporal reduction ayuda.

\lecturefigure{slide-44-results-table.jpg}{Resultados de audio understanding: baselines convolucionales, Transformer y variantes con pooling/multi-scale.}{Video oficial de Stanford Online 00:44:02--00:44:35; Audio Transformers v1, tabla 1.}

\begin{table}[H]
\centering
\small
\caption{FSD50K classroom-era comparison (mAP; el número de parámetros proviene de la source table).}
\begin{tabular}{lrr}
\toprule
Modelo & mAP & Parameters \\
\midrule
CRNN & 0.417 & 0.96M \\
VGG-like & 0.434 & 0.27M \\
ResNet-18 & 0.373 & 11.3M \\
DenseNet-121 & 0.425 & 12.5M \\
Small Transformer & 0.469 & 0.9M \\
Large 6-layer Transformer & 0.525 & 2.3M \\
Large 6-layer Transformer with pooling & 0.537 & 2.3M \\
\bottomrule
\end{tabular}
\end{table}

\begin{warningbox}{La conclusión razonable de «Adieu Convolutions»}
La tabla muestra que, con estos datos, esta training recipe y este rango de parámetros, las Transformer variants alcanzan un mAP más alto; no controla todos los convolutional baselines modernos ni compara streaming latency, energy, robustness o preentrenamiento con más datos. El título es una declaración sobre una dirección de investigación, no un teorema general.
\end{warningbox}

\teachervoice{Al ponente le sorprendió la mejora que aportan pooling/wavelet, porque añaden muy pocos learned parameters. Debe entenderse como inductive-bias evidence: cambiar la jerarquía temporal en sí puede importar más que simplemente añadir parámetros.}

\subsection{Resumen de la sección}

Raw-waveform understanding no equivale a conectar globalmente todos los samples entre sí. El learned front end primero forma tokens, y después pooling/wavelet reducen la resolución temporal. Los resultados respaldan el uso de Transformer y del multi-scale bias en FSD50K, pero la conclusión sigue limitada por los datos, las métricas y el alcance de los baselines.
